\subsection{Void-Specific Redshift Drift}
  \label{subsec:redshift_drift}

  The effective saturation mechanism also affects the temporal evolution of
  redshift for sources embedded in low-density environments.

  In the standard cosmological framework, the redshift drift depends solely on the
  global expansion history.
  In the present framework, sources located in deep voids experience a modified
  effective transition rate, leading to a small but systematic deviation in the
  expected redshift drift signal.

  Future high-precision spectroscopic experiments may be able to detect this
  effect by comparing sources located in voids with those in denser regions,
  complementing existing redshift-drift proposals~\cite{Cai2017VoidLensing}.
  The effect is predicted to be suppressed at high redshift, where environmental
  averaging dominates.
